\documentclass[11pt,a4paper]{article}

\usepackage[T1]{fontenc}
\usepackage{lmodern}
\usepackage{textcomp}
\usepackage{amsmath,amssymb,amsthm}
\usepackage{mathtools}
\usepackage{graphicx}
\usepackage{array}
\usepackage{geometry}
\usepackage[hidelinks]{hyperref}
\usepackage{doi}
\geometry{margin=2.5cm}
\setlength{\emergencystretch}{2em}
\hypersetup{
  pdftitle={Finite Criteria for Simplicial States and Local Tomography toward Emergent Operational Composition},
  pdfauthor={Jérôme Beau},
  pdfsubject={Convexity, effect availability, and local tomography for a finite rank-one admissibility contract},
  pdfkeywords={convex state spaces, generalized probabilistic theories, effect algebras, local tomography, classical randomisation, simplicial geometry, non-injective admissibility, emergent composition}
}

\input{glyphtounicode}
\pdfgentounicode=1

\graphicspath{{out/}{../out/}}

\newtheorem{theorem}{Theorem}[section]
\newtheorem{lemma}[theorem]{Lemma}
\newtheorem{corollary}[theorem]{Corollary}
\newtheorem{proposition}[theorem]{Proposition}
\newtheorem{definition}[theorem]{Definition}
\newtheorem{remark}[theorem]{Remark}
\newtheorem{example}[theorem]{Example}
\newtheorem{axiom}[theorem]{Axiom}

\newcommand{\SA}{S_A}
\newcommand{\SB}{S_B}
\newcommand{\FA}{F_A}
\newcommand{\FB}{F_B}
\newcommand{\FAB}{F_{AB}}
\newcommand{\conv}{\operatorname{conv}}
\newcommand{\rk}{\operatorname{rank}}
\newcommand{\dd}{\delta}
\newcommand{\abs}[1]{\lvert #1 \rvert}

\title{Finite Criteria for Simplicial States and Local Tomography\\
toward Emergent Operational Composition}

\author{J\'{e}r\^{o}me Beau\\
\small Independent Researcher}

\date{2026\\[2pt]\small Version 1.0}

\begin{document}

  \maketitle

  \begin{abstract}
    The companion paper on effective composition characterises, for a finite substrate without a supplied
    subsystem factorisation, the possibilistic support layer completely and the preparation layer for two
    reference classes, and explicitly leaves open the states-and-effects layer of the resulting composition
    contract.
    We resolve that layer's finite, mono-system, and bipartite-composition questions, keeping the general
    preparation layer and the tensor product exactly as open as the companion paper left them.
    The primitive family arising from cylindrical conditioning of a strictly positive, rank-one reference
    measure is not itself convex: it is finite while its own convex hull is a continuum, unconditionally, for
    any local outcome set of size at least two.
    An explicit classical-randomisation postulate closes it onto the full classical probability simplex,
    independently of the reference measure's rank; rank governs composition of the primitive family, not the
    shape of its closure once randomisation is added.
    Simplicial geometry fixes the mathematical representation of affine effects as vectors in a cube, but not
    their physical availability as measurements; a weak coarse-graining axiom already yields the entire
    Boolean skeleton of sharp effects for free, while the full continuum of effects requires a further,
    dual randomisation postulate on the measurement side.
    Under two further, distinct, and logically independent axioms --- one for a common classical control
    resource, one for the physical existence of a composed bipartite effect --- the pointwise product of two
    local effects is proved unique, not postulated, and local tomography holds exactly when the local effect
    families linearly span their ambient functional spaces, equivalently when they separate their local
    states and contain the unit effect in their linear span.
    All results are finite, exact, and reproduced by committed scripts using rational arithmetic.
    \emph{Interpretive reading (structural, not a result):} the finite hypotheses examined here force,
    successively and for free, a classical, simplicial operational skeleton --- convex states once mixing is
    supplied, a Boolean effect algebra once coarse-graining is supplied, and exact local tomography once
    spanning local families are supplied --- under precisely identified imports, and produce neither a
    quantum state geometry nor the Born rule.
    Locating exactly how much operational structure a finite admissibility contract yields for free, and
    where every further richness must instead be imported, is the intended service of this result to the
    non-injective foundations programme: it marks the boundary this classical skeleton cannot cross on its
    own, and hands that boundary, precisely typed, to whatever module or phase law is proposed next.
  \end{abstract}

  \paragraph{Keywords.}
    convex state spaces; generalized probabilistic theories; effect algebras; local tomography; classical
    randomisation; simplicial geometry; non-injective admissibility; emergent composition.

    \clearpage
    \tableofcontents

    \clearpage
  \section{Introduction}
  \label{sec:intro}

  \subsection{Dependency on Effective Composition}
    \label{ssec:dependency}

    The companion paper~\cite{Beau2026esub} studies a finite configuration set $\Omega$ carrying a reference
    admissible subset and two effective projections $\Pi_A, \Pi_B$ that are not assumed injective, with no
    supplied subsystem factorisation $\Omega \simeq \Omega_A \times \Omega_B$.
    It names the missing composition structure a transversal contract,
    $\mathrm{K4} \colon \mathsf{Admissibility} \to \mathsf{Composable\ effective\ subsystems}$, and proves,
    completely at the possibilistic (support) layer and for two reference preparation classes: the
    independence blocks of the projected support relation are exactly its nonempty rectangles (bicliques);
    on a block, the full simplex of preparations lifts every joint law; and, for the family of normalised
    cylindrical conditionings of a strictly positive fibre-weight matrix $W$ on a block, closure under
    products of independently locally conditioned states holds if and only if $\rk W = 1$.
    It explicitly leaves open the general preparation layer and \emph{the states-and-effects layer: the
    construction of effect spaces, their bilinear pairing with states, local tomography, and any quantum
    tensor product}.

    This paper imports the rank-one cylindrical-conditioning family exactly as
    \cite{Beau2026esub}'s Theorem~3 supplies it --- as an already-proved hypothesis on the primitive family,
    not reopened --- and resolves the states-and-effects layer's finite, mono-system, and
    bipartite-composition content.
    The general preparation layer, the tensor product as a supplied structure, local tomography's relation
    to the Born rule, and any phase/modulus law remain exactly as open as the companion paper left them.

  \subsection{Contributions}
    \label{ssec:contributions}

    All results are finite and exact, reproduced by six deterministic scripts using rational arithmetic
    (Appendix~\ref{app:repro}).

    \begin{enumerate}
      \item \textbf{Primitive convexity obstruction (Theorem~\ref{thm:cv1}).}
      The raw cylindrical-conditioning family is finite while its convex hull is a continuum, for any local
      outcome set of size at least two, with no genericity hypothesis --- and the obvious rescue attempt,
      conditioning inside a projection fibre, fails (Remark~\ref{rem:cv3}).
      \item \textbf{Simplicial closure (Theorem~\ref{thm:cv6}).}
      An explicit classical-randomisation postulate closes the family onto the full probability simplex,
      locally and jointly, independently of the reference measure's rank.
      \item \textbf{Representation vs.\ availability of effects (Theorem~\ref{thm:ea1}, Theorem~\ref{thm:ea4}).}
      Affine effects on a simplex correspond to a cube of vectors, a bookkeeping fact special to simplices;
      a weak coarse-graining axiom already yields the full Boolean effect algebra; the full continuum
      requires a further, dual randomisation postulate, realised by an explicit multilinear reconstruction.
      \item \textbf{Independence and consistency of the two randomisations (Theorem~\ref{thm:ccu5}).}
      The state-side and effect-side randomisation postulates are logically independent, and a common
      classical control resource reproduces both only under an explicit independence clause between its two
      interfaces --- dropping it breaks bilinear consistency, by an exact numeric witness.
      \item \textbf{Uniqueness of the composed effect (Theorem~\ref{thm:bc1}).}
      Given physical existence of a bipartite effect and the natural compatibility with joint Dirac states,
      it is forced to equal the pointwise-product formula: no tensor product is imported to define it.
      \item \textbf{Inaccessible joint effects (Theorem~\ref{thm:bc4}).}
      Local composition reaches only rectangle effects, strictly fewer than all joint sharp effects; an
      explicit separating witness shows this survives full local randomisation.
      \item \textbf{Exact criterion for local tomography (Theorem~\ref{thm:lt3}, Corollary~\ref{cor:lt3corollary}).}
      Local tomography holds exactly when each side's effect family linearly spans its ambient space,
      strictly stronger than merely separating its own simplex (Example~\ref{ex:lt2}); equivalently, when it
      separates states and its linear span contains the unit effect.
    \end{enumerate}

  \subsection{What this paper does not claim}
    \label{ssec:nonclaims}

    This paper derives no Born rule, no quantum state geometry, no phase or interference structure, and no
    tensor product of state spaces as a derived object: the pointwise-product formula for effects is forced
    given physical existence of a composed effect (an explicit axiom), not a substitute derivation of tensor
    products for states.
    Every closure beyond the forced discrete/Boolean skeleton --- state convexity, effect convexity, their
    joint consistency, and full joint correlational structure --- is obtained only under an explicitly typed,
    separately motivated postulate, never from the admissibility contract or the simplicial geometry alone.
    The general preparation layer of~\cite{Beau2026esub} remains open; this paper works entirely within the
    strictly positive, rank-one, finite reference class it supplies.

  \section{Primitive Convexity Obstruction}
  \label{sec:cv}

    Fix finite outcome sets $\SA$, $\SB$ and a strictly positive rank-one reference measure
    $W = (w_{ab}) = \alpha_a \beta_b$ on $\SA \times \SB$, imported unchanged from~\cite{Beau2026esub}.
    For nonempty $U \subseteq \SA$, the cylindrical conditioning is
    $p_U^A := \alpha|_U / \alpha(U)$, and symmetrically $p_V^B$ on $\SB$.

    \begin{definition}[Raw families]
      \label{def:raw}
      $\FA := \{\, p_U^A : \emptyset \neq U \subseteq \SA \,\}$, $\FB$ symmetric, and
      $\FAB := \{\, p_U^A \otimes p_V^B : U, V \,\}$.
    \end{definition}

    \begin{theorem}[Primitive family is never itself convex]
      \label{thm:cv1}
      For $m := \abs{\SA} \geq 2$ and any strictly positive $\alpha$: $\FA$ has at most $2^m - 1$ elements
      (finite), $\conv(\FA) = \Delta(\SA)$ (a continuum of dimension $m - 1$), and hence
      $\FA \subsetneq \conv(\FA)$, with no genericity hypothesis on $\alpha$.
    \end{theorem}

    \begin{proof}
      $\abs{\FA} \leq 2^m - 1 < \infty$ directly from Definition~\ref{def:raw}.
      Since every singleton $\{a\}$ is an admissible cylindrical condition, $p_{\{a\}}^A$ is the Dirac vertex
      $\dd_a$, so $\FA$ contains every vertex of $\Delta(\SA)$; hence $\conv(\FA) = \Delta(\SA)$, which has
      dimension $m - 1 \geq 1$ and is therefore uncountable.
      A finite set cannot equal an uncountable one.
    \end{proof}

    \begin{corollary}[Joint family, same obstruction, wider gap]
      \label{cor:cv2}
      $\abs{\FAB} \leq (2^m - 1)(2^n - 1)$ (finite, $n := \abs{\SB}$), while $\Delta(\SA \times \SB)$ has
      dimension $mn - 1$; the already-proved product-closure of~\cite{Beau2026esub} (equivalent to
      $\rk W = 1$) governs composition of $\FAB$ and does not rescue its convex closure --- the two closure
      properties are logically independent.
    \end{corollary}

    \begin{remark}[Objection and reply: sub-fibre conditioning does not rescue convexity]
      \label{rem:cv3}
      Since the underlying projections are non-injective, one might hope that conditioning inside a fibre,
      rather than only on saturated unions of fibres, fills in the continuum for free.
      It does not: a local condition, by construction, is a function of $\Pi_A$-observable information
      alone, and points within one fibre are by definition indistinguishable under $\Pi_A$.
      Any rule that nonetheless selects a sub-fibre slice must consult information external to what $\Pi_A$
      can see --- exactly an external classical selector of the kind whose necessity is in question, now
      relocated rather than removed.
    \end{remark}

  \section{Simplicial Closure under Classical Randomisation}
  \label{sec:closure}

    \begin{axiom}[P1: classical randomisation of preparations]
      \label{ax:p1}
      If $p, q$ are effective preparations already reached and $\lambda \in [0,1]$, an external procedure
      exists that prepares $\lambda p + (1-\lambda) q$ as an effective preparation.
    \end{axiom}

    P1 is imported by name, not derived, and is exactly the ingredient Theorem~\ref{thm:cv1} shows missing.

    \begin{theorem}[Full simplicial closure, rank-independent]
      \label{thm:cv6}
      Under Axiom~\ref{ax:p1}, $\conv(\FA) = \Delta(\SA)$ and $\conv(\FAB) = \Delta(\SA \times \SB)$ exactly
      --- and the joint statement holds for \emph{any} strictly positive $W$, not only $\rk W = 1$.
    \end{theorem}

    \begin{proof}
      The local statement restates Theorem~\ref{thm:cv1}'s vertex argument under P1.
      For the joint statement: $p_{\{a\}}^A \otimes p_{\{b\}}^B = \dd_a \otimes \dd_b = \dd_{(a,b)}$ for every
      $(a,b)$, using only that $p_U^A, p_V^B$ condition the two \emph{marginals} of $W$ --- true whether or
      not $W$ factorises.
      So $\FAB$ contains every joint Dirac vertex regardless of $\rk W$, and P1 reaches their convex hull,
      the whole simplex.
    \end{proof}

    \begin{remark}[Rank governs the generators, not the closure]
      \label{rem:cv7}
      $\rk W$ decides whether the \emph{primitive} family composes as a product (\cite{Beau2026esub}'s
      Theorem~3); it becomes invisible once P1 is added, since a rank-one and a fully correlated $W$ both
      close, under P1, onto the identical $\Delta(\SA \times \SB)$.
    \end{remark}

  \section{Affine Representation vs.\ Physical Availability of Effects}
  \label{sec:ea}

    \begin{definition}[Effect]
      \label{def:effect}
      An effect on $\Delta(S)$ is an affine map $e \colon \Delta(S) \to [0,1]$.
    \end{definition}

    \begin{theorem}[Representation]
      \label{thm:ea1}
      The set of affine $[0,1]$-valued functions on $\Delta(S)$, $\abs{S} = m$, corresponds bijectively to
      $[0,1]^S$, via $v_a := e(\dd_a)$, $e(p) = \sum_a p(a) v_a$.
      This identifies a purely mathematical function space with a cube and carries no physical content: it
      is never, by itself, a claim that every vector of that cube is a physically available effect --- which
      of them are is the separate question this section and Section~\ref{sec:composition} answer.
    \end{theorem}

    \begin{proof}
      Every $p \in \Delta(S)$ decomposes uniquely as $p = \sum_a p(a) \dd_a$, since the Dirac vertices are
      affinely independent (the defining property of a simplex); affinity of $e$ then gives
      $e(p) = \sum_a p(a) e(\dd_a)$.
      Conversely, for any $v \in [0,1]^S$, $p \mapsto \langle v, p \rangle$ is affine and, being a convex
      combination of numbers in $[0,1]$, valued in $[0,1]$.
      The two constructions are mutually inverse.
    \end{proof}

    \begin{remark}[Why this is special to a simplex]
      \label{rem:ea1-simplex}
      Both directions used affine independence of the vertices: it is what makes the decomposition of $p$
      unique, and what removes any consistency constraint on an arbitrary choice of boundary values $v_a$.
      On a non-simplicial convex body, the vertices satisfy nontrivial affine relations, and not every
      boundary assignment extends to an affine function; the representation theorem's triviality here is a
      structural feature of the classical/simplicial case, not a general fact about convex state spaces.
    \end{remark}

    \begin{theorem}[Tight minimal separating family]
      \label{thm:ea2}
      $\{ e_a : a \in S \setminus \{a_0\} \}$ (any $m-1$ coordinate effects) separates $\Delta(S)$; no $m-2$
      of them do.
    \end{theorem}

    \begin{proof}
      If $p, q$ agree on every coordinate but $a_0$, normalisation forces
      $p(a_0) = 1 - \sum_{a \neq a_0} p(a) = 1 - \sum_{a \neq a_0} q(a) = q(a_0)$, so $p = q$.
      For tightness, exhibit $p, q$ agreeing on $m-2$ fixed coordinates while differing on the remaining two,
      both summing to~$1$.
    \end{proof}

    \begin{axiom}[AX-CG: coarse-graining]
      \label{ax:cg}
      Grouping the outcomes of an already-given partition into a coarser yes/no question is itself a
      legitimate effect.
    \end{axiom}

    \begin{theorem}[Boolean skeleton, unit effect for free]
      \label{thm:ea3}
      Under Axiom~\ref{ax:cg} alone, the canonical partition $\{e_a : a \in S\}$ generates exactly the
      $2^m$ indicator effects $\{\mathbf{1}_U : U \subseteq S\}$ --- the full vertex set of $[0,1]^S$ --- and
      its own sum is the unit effect: $\sum_a e_a = u$, with no separate normalisation postulate.
    \end{theorem}

    \begin{proof}
      $\mathbf{1}_U = \sum_{a \in U} e_a$ by coarse-graining; ranging $U$ over all subsets gives exactly the
      $\{0,1\}$-valued vectors, which are exactly the vertices of the cube.
      $\sum_a e_a(p) = \sum_a p(a) = 1$ for every $p$.
    \end{proof}

    \begin{axiom}[AX-RND: effect randomisation, dual of P1]
      \label{ax:rnd}
      If $e_1, e_2$ are effects and $\lambda \in [0,1]$, an external procedure (flip a $\lambda$-biased coin,
      then run $e_1$ or $e_2$) realises $\lambda e_1 + (1-\lambda) e_2$ as an effect.
    \end{axiom}

    \begin{theorem}[Full cube, explicit reconstruction]
      \label{thm:ea4}
      Under Axioms~\ref{ax:cg} and~\ref{ax:rnd}, $\conv(\{\mathbf{1}_U\}) = [0,1]^S$: every
      $x \in [0,1]^S$ equals $\sum_{U \subseteq S} w_U \mathbf{1}_U$ with
      $w_U := \prod_{a \in U} x_a \prod_{a \notin U} (1-x_a)$, the probabilities of $m$ independent Bernoulli
      trials with biases $x_a$.
    \end{theorem}

    \begin{proof}
      $w_U \geq 0$ since each $x_a \in [0,1]$, and $\sum_U w_U = \prod_a (x_a + (1-x_a)) = 1$ by the binomial
      expansion; substitution confirms the sum reproduces $x$ coordinatewise.
    \end{proof}

    \begin{remark}[Topological closure is redundant here]
      \label{rem:ea5}
      $S$ is finite, so any convex hull of the resulting finite vertex set is automatically closed; a
      further topological-closure axiom has no work to do at this scale.
    \end{remark}

  \section{Independence of the Two Randomisations}
  \label{sec:cc}

    P1 (Axiom~\ref{ax:p1}) randomises preparations upstream; AX-RND (Axiom~\ref{ax:rnd}) randomises effects
    downstream.
    Formal duality does not imply operational equivalence: a model may grant either without the other.

    \begin{example}[Independence counter-models]
      \label{ex:cc-indep}
      $\Delta(S_A)$ full, but effects restricted to the $2^m$ Boolean indicators (P1 without AX-RND): every
      pairing $\mathbf{1}_U(p)$ is well defined for $p \in \Delta(S_A)$ regardless of what else is available.
      Dually, the raw unclosed family $\FA$ as states but the full cube $[0,1]^{S_A}$ as effects (AX-RND
      without P1): Theorem~\ref{thm:ea4}'s reconstruction never used state-space convexity.
      Neither postulate presupposes the other.
    \end{example}

    \begin{axiom}[CC: common classical control resource]
      \label{ax:cc}
      An external classical random source, statistically independent of the system, can be wired (CC-P) to
      select between two preparation procedures and, independently (CC-M), to select between two measurement
      procedures; when both interfaces are exercised in the same trial, the two samples drawn are
      statistically independent of each other (CC-IND).
    \end{axiom}

    CC-P and CC-M each immediately deliver P1 and AX-RND respectively, by the law of total probability; the
    substantive content is CC-IND.

    \begin{theorem}[Bilinear consistency requires CC-IND]
      \label{thm:ccu5}
      The algebraic identity $(\lambda e + (1-\lambda)f)(\mu p + (1-\mu)q) = \lambda\mu\, e(p) +
      \lambda(1-\mu)\, e(q) + (1-\lambda)\mu\, f(p) + (1-\lambda)(1-\mu)\, f(q)$ holds unconditionally.
      Physically running a $\lambda$-biased choice of effect and a $\mu$-biased choice of preparation
      reproduces this value \emph{if and only if} the two choices are drawn independently.
      Reusing one shared bit for both choices (forcing $\mu = \lambda$) instead gives
      $\lambda\, e(p) + (1-\lambda)\, f(q)$, missing both cross terms.
    \end{theorem}

    \begin{proof}
      The algebraic identity is direct substitution.
      Independent draws average the reported outcome over four branches with product weights
      $\lambda\mu, \lambda(1-\mu), (1-\lambda)\mu, (1-\lambda)(1-\mu)$, reproducing the identity by
      construction.
      A single shared bit has only two branches, weighted $\lambda, 1-\lambda$, giving the stated two-term
      sum.
    \end{proof}

    \begin{example}[Numeric witness]
      \label{ex:ccu5-numeric}
      On $S = \{0,1\}$, with $p = (7/10, 3/10)$, $q = (1/5, 4/5)$, $e = (9/10, 1/10)$, $f = (2/5, 3/5)$,
      $\lambda = 1/3$: independent taps give exactly $217/450$, matching the algebraic value; a shared bit
      at $\mu = \lambda = 1/3$ gives $89/150 = 267/450$, a discrepancy of exactly $1/9$
      (\texttt{code/cc\_unification\_tests.py}).
    \end{example}

    \begin{remark}[Verdict]
      \label{rem:ccu6}
      Unification of P1 and AX-RND under a single resource is possible, but only under Axiom~\ref{ax:cc} in
      full, including CC-IND; it is not derived from $\Delta(S)$'s geometry, nor from P1/AX-RND's formal
      duality alone.
    \end{remark}

  \section{Bipartite Composition: the Pointwise Product is Forced}
  \label{sec:composition}

    \begin{axiom}[AX-COMP: physical existence of a composed effect]
      \label{ax:comp}
      For each pair of physically available local effects $e$ (on $\SA$), $f$ (on $\SB$), a joint effect
      $e \boxtimes f$, affine on $\Delta(\SA \times \SB)$, is physically realised.
    \end{axiom}

    \begin{theorem}[Uniqueness of the composed effect]
      \label{thm:bc1}
      If $e \boxtimes f$ exists (Axiom~\ref{ax:comp}) and satisfies Dirac-factorisation,
      $(e \boxtimes f)(\dd_{(a,b)}) = e(\dd_a) f(\dd_b)$, then it is uniquely
      \begin{equation}
        \label{eq:pointwise}
        (e \boxtimes f)(p) = \sum_{(a,b)} p_{ab}\, e_a f_b.
      \end{equation}
    \end{theorem}

    \begin{proof}
      Immediate corollary of Theorem~\ref{thm:ea1} applied with outcome set $\SA \times \SB$:
      $\Delta(\SA \times \SB)$ is itself a simplex, so any affine function on it (which $e \boxtimes f$ is,
      by Definition~\ref{def:effect}) is fixed uniquely by its Dirac values, given as $e_a f_b$ by
      hypothesis.
    \end{proof}

    Equation~\eqref{eq:pointwise} is what one would ordinarily call the tensor-product effect $e \otimes f$
    evaluated pointwise; no tensor-product formalism is imported to \emph{define} it.
    What remains imported is Axiom~\ref{ax:comp}'s existence clause: for a given pair, whether such a joint
    measurement is physically available at all.

    \begin{remark}[Operational bilinearity again requires CC-IND]
      \label{rem:bc2}
      As algebra, $e \boxtimes f$ is bilinear in $(e,f)$ with no ambiguity.
      The operational question --- does physically randomising the local-effect choice on each side
      reproduce the bilinear expansion in a single trial --- is exactly Theorem~\ref{thm:ccu5}, relocated to
      the two composition interfaces.
      Numeric witness (\texttt{code/bipartite\_composition\_tests.py}, generic non-product joint state):
      independent taps reproduce the bilinear value $161/600$ exactly; a shared bit at the forced common
      bias gives $147/500$ against a true value of $49/180$, a discrepancy of $49/2250$.
    \end{remark}

  \section{Inaccessible Joint Effects: the Diagonal Witness}
  \label{sec:diagonal}

    Taking $e = \mathbf{1}_U$, $f = \mathbf{1}_V$ (local sharp effects), Equation~\eqref{eq:pointwise} gives
    $\mathbf{1}_U \boxtimes \mathbf{1}_V = \mathbf{1}_{U \times V}$, the indicator of the rectangle
    $U \times V$.

    \begin{theorem}[Rectangles are strictly fewer than all joint sharp effects]
      \label{thm:bc3}
      For $m, n \geq 2$, rectangles number at most $(2^m-1)(2^n-1) + 1$, strictly fewer than $2^{mn}$; for
      $m = n = 2$, exactly $10$ of the $16$ subsets of $\SA \times \SB$ are rectangles.
    \end{theorem}

    \begin{proof}[Proof (census, $m=n=2$)]
      By size: $1$ (empty) $+ 4$ (singletons, all rectangles) $+ 4$ (the four row/column pairs, out of $6$
      possible size-two subsets) $+ 0$ (no size-three subset has $\abs{U}\abs{V} = 3$ for
      $\abs{U},\abs{V} \in \{1,2\}$) $+ 1$ (the whole space) $= 10$
      (\texttt{code/bipartite\_composition\_tests.py}).
    \end{proof}

    \begin{theorem}[Local products, however randomised, cannot reach the diagonal]
      \label{thm:bc4}
      Define $\varphi(M) := M_{00} + M_{11} - M_{01} - M_{10}$ on $\SA \times \SB = \{0,1\}^2$.
      For any convex combination of rank-one products $M = \sum_k w_k(e_k \otimes f_k)$,
      $\varphi(M) \in [-1,1]$; but $\varphi(\mathbf{1}_{\{(0,0),(1,1)\}}) = 2$.
    \end{theorem}

    \begin{proof}
      $\varphi(e \otimes f) = (e_0 - e_1)(f_0 - f_1)$, a product of two numbers in $[-1,1]$ (since
      $e_i, f_i \in [0,1]$), hence itself in $[-1,1]$; $\varphi$ is linear, so a weighted average of such
      values stays in $[-1,1]$.
      The diagonal indicator has $\varphi = 1 + 1 - 0 - 0 = 2$, outside that range
      (\texttt{code/bipartite\_composition\_tests.py}).
    \end{proof}

    \begin{remark}[A structural resemblance, not a physics claim]
      \label{rem:chsh}
      The bound $\varphi \in [-1,1]$ for mixtures of product objects, violated by a genuinely correlated one,
      has the shape of a CHSH-type inequality~\cite{ClauserHorneShimonyHolt1969}.
      This is stated here as an elementary fact about real numbers in $[-1,1]$ and convex combinations; it
      makes no claim about, and is not evidence for or against, anything in the physical Bell-nonlocality
      literature, which concerns quantum measurement statistics, a different question.
    \end{remark}

    \begin{remark}[Headline classification]
      \label{rem:bc5}
      Reaching the full joint Boolean algebra (all $2^{mn}$ sharp effects) or the full cube
      $[0,1]^{\SA \times \SB}$ requires a separate, strictly stronger axiom applied directly to
      $\SA \times \SB$ as its own system (Axioms~\ref{ax:cg}, \ref{ax:rnd} run with that combined outcome
      set) --- not obtainable from local composition, local randomisation, or their combination.
    \end{remark}

  \section{Exact Criterion for Local Tomography}
  \label{sec:tomography}

    \begin{theorem}[Full coordinate families reconstruct the joint state]
      \label{thm:lt1}
      $(e_a \boxtimes f_b)(p) = p_{ab}$ for every $(a,b)$ (immediate from Equation~\eqref{eq:pointwise}); the
      $mn$ numbers so obtained are $p$'s own coordinates.
    \end{theorem}

    \begin{example}[The naive general formulation is false]
      \label{ex:lt2}
      Take $E_A = \{e_1\}$, $E_B = \{f_1\}$ on $\SA = \SB = \{0,1\}$ --- each separates its own simplex
      (Theorem~\ref{thm:ea2}, $m=2$).
      Yet $q = (1/10, 2/10, 3/10, 4/10)$ and $q' = (3/10, 1/10, 2/10, 4/10)$ (order $(0,0),(0,1),(1,0),(1,1)$)
      are distinct states with $(e_1 \boxtimes f_1)(q) = (e_1 \boxtimes f_1)(q') = 2/5$
      (\texttt{code/local\_tomography\_tests.py}).
      Local separation of each side does not imply that products separate the joint simplex.
    \end{example}

    \begin{theorem}[Corrected criterion: linear spanning]
      \label{thm:lt3}
      If $E_A$'s vectors linearly span $\mathbb{R}^{\SA}$ and $E_B$'s span $\mathbb{R}^{\SB}$, then
      $E_A \boxtimes E_B$ separates $\Delta(\SA \times \SB)$.
    \end{theorem}

    \begin{proof}
      Let $D := Q - Q'$ satisfy $e^T D f = 0$ for all $e \in E_A, f \in E_B$.
      Fixing $e$, the row vector $e^T D$ is orthogonal to every $f \in E_B$; since $E_B$ spans
      $\mathbb{R}^{\SB}$, $e^T D = 0$ for every $e \in E_A$.
      Since $E_A$ spans $\mathbb{R}^{\SA}$, $D = 0$.
      Contrapositive: $Q \neq Q'$ forces some $e,f$ with $(e \boxtimes f)(Q) \neq (e \boxtimes f)(Q')$.
    \end{proof}

    Spanning needs at least $m$ elements, strictly more than Theorem~\ref{thm:ea2}'s $m-1$-sufficient
    separation; Example~\ref{ex:lt2} is exactly the gap between the two.

    \begin{corollary}[Separation plus the unit effect is equivalent to spanning]
      \label{cor:lt3corollary}
      If $E$ separates $\Delta(S)$ and the unit effect $u = (1,\dots,1) \in \operatorname{span}(E)$, then
      $\operatorname{span}(E) = \mathbb{R}^S$.
    \end{corollary}

    \begin{proof}
      Suppose $V := \operatorname{span}(E) \subsetneq \mathbb{R}^S$ with $u \in V$.
      Take nonzero $w$ in $V^\perp$: $w \perp e$ for every $e \in E$, and $w \perp u$, i.e.\
      $\sum_a w_a = 0$.
      Fix an interior $p_0 \in \Delta(S)$; for small enough $\varepsilon$, $p := p_0 + \varepsilon w$ and
      $p' := p_0 - \varepsilon w$ both lie in $\Delta(S)$ (their coordinates still sum to $1$) and are
      distinct.
      Yet $e(p) - e(p') = 2\varepsilon \langle e, w \rangle = 0$ for every $e \in E \subseteq V$: $E$ fails to
      separate $p$ from $p'$, contradicting the hypothesis.
      So no such proper $V$ exists.
    \end{proof}

    Theorem~\ref{thm:ea2}'s minimal family spans only the coordinate hyperplane $\{x : x_{a_0} = 0\}$, which
    does not contain $u$; adding $u$ supplies precisely the missing direction, and under Axiom~\ref{ax:cg}'s
    full partition $u$ is already free (Theorem~\ref{thm:ea3}).
    Local tomography therefore needs nothing beyond a locally separating family, composable by
    Axiom~\ref{ax:comp}, enriched by its own normalisation.

    \begin{remark}[Redundancy and a computable-but-unavailable expectation]
      \label{rem:lt5lt6}
      Dropping any single one of the $mn$ full-coordinate products still permits exact reconstruction by
      normalisation, echoing Theorem~\ref{thm:ea2} at the joint level.
      Once tomography reconstructs $p$, the expectation of \emph{any} affine functional --- including the
      diagonal indicator of Theorem~\ref{thm:bc4}, physically unavailable by composition alone --- is simply
      $\sum_{ab} p_{ab} v_{ab}$: physical availability of an effect and computability of its expectation on a
      reconstructed state are different properties (\texttt{code/local\_tomography\_tests.py}).
    \end{remark}

  \section{Position and Scope}
  \label{sec:position}

  \subsection{Position within the neighbouring literature}
    \label{ssec:position}

    \paragraph{Generalized probabilistic theories.}
    States, effects, and their bilinear pairing are standard operational input in GPT
    reconstructions~\cite{Hardy2001,Barrett2007,ChiribellaDArianoPerinotti2011}; local tomography is among
    the axioms distinguishing complex quantum theory from its real and quaternionic
    variants~\cite{Barrett2007}, and recent work shows single-system structure alone does not fix the
    composition rule~\cite{ErbaPerinotti2024}.
    This paper works one level below that literature's starting point: it derives which parts of the
    state-effect-tomography package a finite admissibility contract supplies for free, and which require an
    explicitly typed import, rather than taking the full package as given.

    \paragraph{Operational effect algebras.}
    Effects as affine $[0,1]$-valued functionals on a convex state space, and their operational role, trace
    to~\cite{DaviesLewis1970} and to Ludwig's axiomatic convex-state-space
    programme~\cite{Ludwig1983}.
    Theorems~\ref{thm:ea1}--\ref{thm:ea4} work inside that same formal object but ask a different question of
    it: which subset of the mathematically possible effect functionals is forced by coarse-graining alone,
    and which requires an explicit randomisation postulate, rather than assuming from the outset that every
    bounded affine functional is available.

    \paragraph{CHSH and Bell-inequality mathematics.}
    The separating witness of Theorem~\ref{thm:bc4} has the algebraic shape of the original
    CHSH construction~\cite{ClauserHorneShimonyHolt1969}, restricted here to an entirely classical,
    finite-dimensional statement about real numbers in $[-1,1]$; Remark~\ref{rem:chsh} states explicitly that
    no claim about physical nonlocality is made or needed.

  \subsection{Scope}
    \label{ssec:scope}

    The finite hypotheses studied here force, successively and under precisely identified imports, a
    classical operational skeleton: convex states once classical randomisation (Axiom~\ref{ax:p1}) is
    supplied; a Boolean effect algebra once coarse-graining (Axiom~\ref{ax:cg}) is supplied, and the full
    effect cube once its dual randomisation (Axiom~\ref{ax:rnd}) is also supplied; a forced, tensor-free
    pointwise product once physical composition (Axiom~\ref{ax:comp}) is supplied; and exact local tomography
    once the local families span their ambient spaces.
    Nowhere in this chain does a quantum state geometry or the Born rule appear, and none is derived: every
    non-classical structure the wider programme eventually needs --- a complex or quaternionic carrier, an
    interference/modulus law, a non-simplicial state space --- lies strictly outside what this contract
    supplies, and must be imported by its own, separately typed postulate.

  \section{Discussion}
  \label{sec:discussion}

  \subsection{Interpretive outlook}
    \label{ssec:outlook}

    The following reading is interpretation --- a structural way of situating the results --- and not itself
    a result.
    A finite admissibility contract, left to itself, does not drift toward quantum structure; it settles, at
    every layer examined, into the classical/simplicial regime, and only moves beyond it when a
    correspondingly explicit postulate is supplied.
    This is not a deficiency of the contract: it is a precise map of where the classical skeleton's authority
    ends, handed forward, typed axiom by typed axiom (P1, AX-CG, AX-RND, CC, AX-COMP), to whatever
    phase/modulus law is proposed to take the programme beyond it.
    A future amplitude-modulus construction inherits this map as its starting boundary condition, not as a
    competing description to be reconciled with it.

  \subsection{Limits and future work}
    \label{ssec:limits}

    The general preparation layer of~\cite{Beau2026esub} is untouched.
    The CHSH-type witness of Theorem~\ref{thm:bc4} is proved for $m=n=2$; whether it generalises cleanly to
    larger outcome sets is not shown here, though the rectangle/all-subsets gap of Theorem~\ref{thm:bc3} only
    widens.
    Whether Axiom~\ref{ax:cc}'s independence clause should itself be independently motivated, rather than
    postulated, is left open.
    The natural next step, outside this contract entirely, is to ask what could constrain an amplitude
    modulus without conflating exclusion, mixture, and coherent cancellation --- a question this paper
    deliberately leaves untouched.

  \section*{Acknowledgements}
    Portions of the analysis and editorial preparation benefited from iterative interactions with large
    language models used as analytical assistants.
    All claims and final formulations remain the sole responsibility of the author.

  \appendix

  \section{Reproduction}
  \label{app:repro}

    All numerical claims are reproduced by six deterministic scripts using exact rational arithmetic
    (Python's \texttt{fractions.Fraction}) and no dependencies beyond the standard library:
    \begin{itemize}
      \item \texttt{code/convexity\_obstruction\_tests.py} --- Theorem~\ref{thm:cv1}, Corollary~\ref{cor:cv2};
      \item \texttt{code/simplicial\_closure\_tests.py} --- Theorem~\ref{thm:cv6};
      \item \texttt{code/affine\_effects\_tests.py} --- Theorems~\ref{thm:ea1}--\ref{thm:ea4};
      \item \texttt{code/cc\_unification\_tests.py} --- Theorem~\ref{thm:ccu5}, Example~\ref{ex:ccu5-numeric};
      \item \texttt{code/bipartite\_composition\_tests.py} --- Theorem~\ref{thm:bc1},
      Remark~\ref{rem:bc2}, Theorems~\ref{thm:bc3}--\ref{thm:bc4};
      \item \texttt{code/local\_tomography\_tests.py} --- Theorem~\ref{thm:lt1}, Example~\ref{ex:lt2},
      Theorem~\ref{thm:lt3}, Corollary~\ref{cor:lt3corollary}, Remark~\ref{rem:lt5lt6}.
    \end{itemize}
    Figure~\ref{fig:ladder} is generated by \texttt{code/make\_figures.py}, under the pinned environment
    recorded in \texttt{code/requirements.txt}.
    From a fresh clone:
    \begin{verbatim}
python3 code/convexity_obstruction_tests.py
python3 code/simplicial_closure_tests.py
python3 code/affine_effects_tests.py
python3 code/cc_unification_tests.py
python3 code/bipartite_composition_tests.py
python3 code/local_tomography_tests.py
bash build.sh
    \end{verbatim}
    The scripts print the full inventories; \texttt{build.sh} regenerates the figure, verifies it against
    \texttt{ARTIFACT\_SHA256SUMS}, and compiles the paper.

    \begin{figure}[ht]
      \centering
      \includegraphics[width=\textwidth]{fig_ladder}
      \caption{A: the $7 = 2^3-1$ points reached by cylindrical conditioning on $\Delta(\SA)$, $m=3$
      (Theorem~\ref{thm:cv1}), against the shaded full simplex reachable only under Axiom~\ref{ax:p1}.
      B: a rectangle (reachable by local composition) against the diagonal (unreachable, Theorem~\ref{thm:bc4})
      on $\SA \times \SB = \{0,1\}^2$.
      Generated by \texttt{code/make\_figures.py}.}
      \label{fig:ladder}
    \end{figure}

  \bibliographystyle{plain}
  \bibliography{cosmochrony-bibliography,external-refs}

\end{document}
